\subsection{An exact constant-trace sum and its leading cancellation}
\label{sec:terminal-modal-signed-green-addendum}

The constant-$U$ part of the five-piece split can be summed exactly.  This
does not prove the entry-profile lemma: the varying-$U$, $M_n$, and endpoint
pieces remain.  It does identify the order-one velocity terms that a valid
remainder estimate must cancel.

Put $\gamma=1-q$, $\theta=2s\pi/m$, $\phi=\theta/2$, $z=e^{i\theta}$, and
\begin{equation*}
 R=\frac{\gamma(1+z)}2,\qquad U_0=\frac{3q^3}{40}.
\end{equation*}
For $L\geq0$, write
\begin{equation*}
 \Sigma_L(w)=\sum_{\ell=1}^Lw^\ell,\qquad
 \Gamma_L(w)=\sum_{\ell=0}^Lw^\ell.
\end{equation*}

\begin{lemma}[Small $U$-trace variation under the open chronology]
\label{lem:terminal-modal-u-variation}
Assume the conclusion of \cref{lem:terminal-modal-admission-chronology}.
For $m\geq5$ and $4\leq n\leq m-1$,
\begin{equation}
 0\leq U_0-U_n\leq\frac{q^3}{150},\qquad
 \sum_{n=4}^{m-2}|U_{n+1}-U_n|\leq\frac{q^3}{150}.
 \label{eq:terminal-modal-u-variation}
\end{equation}
\end{lemma}

\begin{proof}
Monotonicity was proved in
\cref{lem:terminal-modal-source-transform}.  To quantify it, write the
right side of \eqref{eq:terminal-modal-frontier-optimum} as
$q^2\mathcal P(t)$, $t=\chi^k$.  On $7/8\leq t\leq1$ and $q\leq1/64$,
the derivative formula from that proof gives
\begin{equation*}
 0\leq\mathcal P'(t)
 =\frac2{1+q}\frac{(1+\chi^{-1})(1-t^2+2t/5)}{(1+t^2)^2}
 \leq\frac{31}{40}.
\end{equation*}
Also $1-\chi^m\leq2/17$: indeed
$2m\operatorname{artanh}q\leq512/4095<\log(17/15)$.
Consequently $q^2\mathcal P(1)-P_{n-1}\leq31q^2/340$.  The value used to upper-bound
$U_n$ at $t=1$ is
\[
 q^3\frac{(1-q)(6/5+2q/5)}{16}
 =q^3\left(\frac3{40}-\frac q{20}-\frac{q^2}{40}\right).
\]
Since $q\eta^2/(4(1+q))\leq q/16$, it follows that
\[
 U_0-U_n\leq q^3\left(
 \frac{31}{5440}+\frac1{1280}+\frac1{163840}\right)
 <\frac{q^3}{150}.
\]
The total-variation claim follows because $U_n$ is nonincreasing.
\end{proof}

\begin{lemma}[Exact constant-$U$ Green sums; conditional on chronology]
\label{lem:terminal-modal-constant-u-sum}
Assume the proper-prefix chronology in
\cref{lem:terminal-modal-admission-chronology}, and let
$X_{m+1}^{U_0}$ and $X_m^{U_0}$ be the constant-$U_0$ components of
$\mathcal A_{m+1}(\theta)$ and $\mathcal A_m(\theta)$ in the five-piece
split.  For $m\geq5$,
\begin{align}
 X_{m+1}^{U_0}
 &=2U_0\Re\left((1+3z)\left[
 z^{-2}\Sigma_{m-4}(Rz^{-2})
 -z^{-1}\Sigma_{m-4}(Rz^{-1})\right]\right),
 \label{eq:terminal-modal-constant-u-position}\\
 X_m^{U_0}
 &=2U_0\Re\left((1+3z)\left[
 z^{-3}\Gamma_{m-5}(Rz^{-2})
 -z^{-2}\Gamma_{m-5}(Rz^{-1})\right]\right).
 \label{eq:terminal-modal-constant-u-previous}
\end{align}
Thus its position and velocity pieces are respectively
$X_{m+1}^{U_0}/q^2$ and
$(X_{m+1}^{U_0}-\gamma X_m^{U_0})/q^3$.

For every fixed $s\geq1$, put $a=1/16$, $b=\pi s$,
$\sigma=(-1)^s$, and $E=e^{-a}$.  Then
\begin{align}
 \lim_{m\to\infty}\frac{X_{m+1}^{U_0}}{q^2}
 &=\frac{3a(\sigma E-1)}{80}
 \left(\frac1{a^2+b^2}-\frac1{a^2+9b^2}\right),
 \label{eq:terminal-modal-constant-u-position-limit}\\
 \lim_{m\to\infty}
 \frac{X_{m+1}^{U_0}-\gamma X_m^{U_0}}{q^3}
 &=\frac{3\sigma b^2(E-\sigma)}5
 \left(\frac1{a^2+b^2}+\frac3{a^2+9b^2}\right).
 \label{eq:terminal-modal-constant-u-velocity-limit}
\end{align}
\end{lemma}

\begin{proof}
The source at prefix $n$ is
\[
 2U_0\Re\left[z^{n-2}(1-z)(1+3z)\right].
\]
Its Green multiplier at time $m+1$ is
\[
 G_{m-n}=R^{m-n}\sum_{j=0}^{m-n}z^{-j}.
\]
Since $(1-z)\sum_{j=0}^Lz^{-j}=z^{-L}-z$ and $z^m=1$, setting
$\ell=m-n$ and summing $4\leq n\leq m-1$ gives
\eqref{eq:terminal-modal-constant-u-position}.  At time $m$ the lag is one
smaller; setting $\ell=m-1-n$ gives
\eqref{eq:terminal-modal-constant-u-previous}.

For the limits, put $x=n/m$.  Uniformly on compact $x$-ranges and with the
endpoints handled by the same bounds,
\begin{align*}
 \frac m{U_0}F_n^{U_0}&\longrightarrow
 16\pi s\sin(2\pi sx),\\
 \frac1mG_{m-n}&\longrightarrow
 e^{-a(1-x)}\frac{\sin(b(1-x))}{b},\\
 G_{m-n}-\gamma G_{m-n-1}&\longrightarrow
 e^{-a(1-x)}\cos(b(1-x)).
\end{align*}
The finite sums are therefore Riemann sums.  Since
$mU_0/q^2=3/640$, their limits are
\begin{align*}
 \frac3{40}\int_0^1e^{-a(1-x)}
 \sin(2bx)\sin(b(1-x))\,dx,\\
 \frac{6b}{5}\int_0^1e^{-a(1-x)}
 \sin(2bx)\cos(b(1-x))\,dx.
\end{align*}
Using $\sin(b(1-x))=(-1)^{s+1}\sin(bx)$, product-to-sum, and
\[
 \int_0^1e^{-a(1-x)}\cos(kbx)\,dx
 =\frac{a(\sigma-E)}{a^2+k^2b^2}\qquad(k=1,3)
\]
gives \eqref{eq:terminal-modal-constant-u-position-limit}.  Similarly,
$\cos(b(1-x))=\sigma\cos(bx)$ and
\[
 \int_0^1e^{-a(1-x)}\sin(kbx)\,dx
 =\frac{kb(E-\sigma)}{a^2+k^2b^2}\qquad(k=1,3)
\]
give
\eqref{eq:terminal-modal-constant-u-velocity-limit}.
\end{proof}

The early base has an equally explicit leading term.  Let $B_{m+1}$ and
$B_m$ denote the homogeneous and early-source components before subtracting
the ideal packet, and put $I_m=mG_{2s}$.  Direct substitution in four literal
prefix steps gives, for fixed $s$,
\begin{equation*}
 \frac{\mathcal A_3}{q^2}=-1+q+O_s(q^2),\qquad
 \frac{\mathcal A_4}{q^2}=-1+\frac75q+O_s(q^2).
\end{equation*}
Solving the subsequent homogeneous scalar recurrence yields
\begin{lemma}[Leading base velocity; conditional on chronology]
\label{lem:terminal-modal-base-velocity-limit}
For every fixed $s\geq1$,
\begin{equation}
 \frac{B_{m+1}-I_m}{q^2}\longrightarrow0,\qquad
 \frac{B_{m+1}-\gamma B_m}{q^3}+\frac{I_m}{5q^2}
 \longrightarrow-\frac45\sigma E.
 \label{eq:terminal-modal-base-velocity-limit}
\end{equation}
\end{lemma}

\begin{proof}
If $r=\gamma\cos\phi$, the homogeneous solution starting at times three and
four is
\[
 B_N=r^{N-3}\left[\mathcal A_3\cos((N-3)\phi)
 +\beta_m\sin((N-3)\phi)\right],
 \quad
 \beta_m=\frac{\mathcal A_4/r-\mathcal A_3\cos\phi}{\sin\phi}.
\]
The displayed prefix expansions give
$\beta_m/q^2\to-3/(80\pi s)$, while
$r^m\to E$, $\cos(m\phi)=\sigma$, and $\sin(m\phi)=0$.  These facts prove
the first limit.  The exact identity
\[
 B_N-\gamma B_{N-1}
 =\gamma r^{N-4}\sin\phi\left[
 \beta_m\cos((N-3)\phi)-\mathcal A_3\sin((N-3)\phi)
 \right]
\]
gives a limit $-3\sigma E/5$ after division by $q^3$.  Finally
$I_m/q^2\to-\sigma E$, proving the second limit.
\end{proof}

There is also an exact reduction for the final piece.  Let
\begin{equation*}
 \mathcal E_{m,s}:=\frac{\widehat F_{m+1}^{\rm ev}(\theta)-A_{m+1}(m)}{q^2}
\end{equation*}
be its position contribution.  Since the newly admitted proper-prefix
coordinate is zero, $e_m(m-1)=-P_m$.  Hence its velocity contribution is
exactly
\begin{equation}
 \frac{\mathcal E_{m,s}}q+\eta\frac{P_m}{q^2}-\frac15.
 \label{eq:terminal-modal-endpoint-velocity-reduction}
\end{equation}
Moreover, with $t=e^{-1/8}$,
\begin{equation}
 \frac{P_m}{q^2}\longrightarrow
 \overline P:=2\frac{2t+(t^2-1)/5}{1+t^2}.
 \label{eq:terminal-modal-frontier-limit}
\end{equation}
Thus the single still-open endpoint estimate
\begin{equation}
 \frac{\mathcal E_{m,s}}q\longrightarrow-\frac15
 \label{eq:terminal-modal-endpoint-open-limit}
\end{equation}
would give the exact final-piece velocity limit
$\overline P/2-2/5$.  Equation
\eqref{eq:terminal-modal-endpoint-open-limit} is measured for fixed modes but
is not proved or used as a theorem.  Together,
\cref{eq:terminal-modal-constant-u-velocity-limit,eq:terminal-modal-base-velocity-limit,eq:terminal-modal-endpoint-velocity-reduction}
make the order-one cancellation explicit; uniform varying-$U$ and $M_n$
control is still required for the velocity profile.
